\begin{abstract}
Spatially continuous mangrove canopy-height estimates are useful for characterizing forest structure, but spaceborne lidar provides sparse observations rather than complete coverage. We investigate sparse-to-dense height regression in Cà Mau, Vietnam, using Sentinel-1 SAR and Sentinel-2 optical composites with sparse GEDI RH98 supervision, without interpolating targets or supplying GEDI to the network at inference. Four U-Net configurations compare optical-only, SAR-only, early-fusion, and dual-encoder feature-fusion strategies under a shared spatial split with buffered boundaries. The reported multi-seed aggregates give the dual-encoder model the lowest test MAE of 2.362 $\pm$ 0.137 m and RMSE of 3.123 $\pm$ 0.179 m, corresponding to approximately 4.7\% and 5.3\% reductions relative to the optical-only baseline. These results suggest that separate modality encoders can improve height prediction in the held-out region, while the small test sample, footprint-to-grid mismatch, and GEDI-only evaluation limit conclusions about accuracy across the dense prediction map.
\end{abstract}
